\begin{figure}[t]
\centering
\begin{tikzpicture}[
    >={Stealth[length=2mm,width=1.6mm]},
    axislab/.style={font=\footnotesize},
    tick/.style={font=\scriptsize},
    fault/.style={dashed, red!70!black, thick},
    flow/.style={thick, blue!60!black},
    jump/.style={thick, red!70!black},
    annot/.style={font=\scriptsize}
  ]

  % Plot occupies x in [0, 8] s.  Top panel translated up.

  % ===== TOP PANEL: survivor cardinality |S(t)| ==============
  \begin{scope}[yshift=2.6cm]
    \draw[->, thick] (-0.2, 0) -- (8.4, 0)
          node[axislab, below right=-1pt] {$t$};
    \draw[->, thick] (0, -0.2) -- (0, 1.85)
          node[axislab, above right=-1pt] {$|\mathcal{S}(t)|$};

    \draw[flow] (0.0, 1.5) -- (2.0, 1.5);
    \draw[flow, dotted] (2.0, 1.5) -- (2.0, 1.0);
    \draw[flow] (2.0, 1.0) -- (5.0, 1.0);
    \draw[flow, dotted] (5.0, 1.0) -- (5.0, 0.5);
    \draw[flow] (5.0, 0.5) -- (8.0, 0.5);

    \node[tick, anchor=east] at (+1.0, 1.7) {$N{=}5$};
    \node[tick, anchor=east] at (+3.0, 1.2) {$N{=}4$};
    \node[tick, anchor=east] at (+6.0, 0.7) {$N{=}3$};

    \draw[fault] (2.0, -0.1) -- (2.0, 1.85);
    \draw[fault] (5.0, -0.1) -- (5.0, 1.85);
    \node[tick, red!70!black, anchor=south] at (2.0, 1.85)
          {$t_1^\star$};
    \node[tick, red!70!black, anchor=south] at (5.0, 1.85)
          {$t_2^\star$};
  \end{scope}

  % ===== BOTTOM PANEL: Lyapunov sawtooth V(t) ================
  \draw[->, thick] (-0.2, 0) -- (8.4, 0)
        node[axislab, below right=-1pt] {$t$};
  \draw[->, thick] (0, -0.2) -- (0, 2.45)
        node[axislab, above right=-1pt]
        {$V(\boldsymbol{\xi}(t))$};

  % flow segment 1
  \draw[flow]
    plot[smooth, samples=40, domain=0:2]
      (\x, {0.4 + 0.9*exp(-0.7*\x)});
  % jump 1
  \draw[jump, ->] (2.0, {0.4 + 0.9*exp(-1.4)})
                -- (2.0, 1.55);
  % flow segment 2
  \draw[flow]
    plot[smooth, samples=40, domain=2:5]
      (\x, {0.4 + 1.15*exp(-0.7*(\x-2))});
  % jump 2
  \draw[jump, ->] (5.0, {0.4 + 1.15*exp(-2.1)})
                -- (5.0, 1.10);
  % flow segment 3
  \draw[flow]
    plot[smooth, samples=40, domain=5:8]
      (\x, {0.4 + 0.7*exp(-0.7*(\x-5))});

  % faults — bottom labels
  \draw[fault] (2.0, -0.1) -- (2.0, 2.45);
  \draw[fault] (5.0, -0.1) -- (5.0, 2.45);
  \node[tick, red!70!black, anchor=north] at (2.0, -0.1)
        {$t_1^\star$};
  \node[tick, red!70!black, anchor=north] at (5.0, -0.1)
        {$t_2^\star$};

  % H2 dwell brace below the bottom panel
  \draw[decorate, decoration={brace, amplitude=2pt},
        thin] (2.0, 4.5) -- (5.0, 4.5)
        node[annot, midway, above=2pt]
        {$\tau_{d,1} \ge \tau_{\mathrm{pend}}$ (H2)};

  % ===== Short on-figure labels =============================
  % Each label is a single word + lemma tag.  Math expressions
  % are in the caption.  Labels live in the upper band, well
  % above the curves (the curves stay below y=1.6).

  % FLOW label — placed above the first decaying segment.
  \node[annot, blue!60!black, anchor=south]
        at (1.0, 1.95)
        {\textbf{flow} (Lem.~\ref{lem:altitude-iss})};

  % JUMP label — placed slightly to the right of jump 1 in
  % white space above the second decaying segment.
  \node[annot, red!70!black, anchor=south]
        at (3.5, 1.85)
        {\textbf{jump} (Lem.~\ref{lem:jump-continuity})};

  % CONTRACTION label — above the third flow segment.
  \node[annot, red!70!black, anchor=south]
        at (6.5, 1.55)
        {\textbf{contraction} (Lem.~\ref{lem:dwell-decay})};

  % chi labels — small white-backed labels left of each jump
  \node[annot, red!70!black, anchor=east, fill=white,
        inner sep=1.2pt] at (1.95, 1.10)
        {$\chi(\Delta_{f,1})$};
  \node[annot, red!70!black, anchor=east, fill=white,
        inner sep=1.2pt] at (4.95, 0.80)
        {$\chi(\Delta_{f,2})$};

\end{tikzpicture}
\caption{Hybrid sawtooth structure of the closed-loop Lyapunov
function across two unannounced cable-severance events.
\emph{Top:} the survivor set $\mathcal{S}(t)$ steps down by one
at each fault $t_k^\star$; cardinality is right-continuous.
\emph{Bottom:} $V(\boldsymbol{\xi})$ exhibits three regimes,
labelled inline. \textbf{Flow:} between events the continuous-time
decay of \cref{lem:altitude-iss} gives
$\dot V \le -\lambda V + \gamma_w(\|w\|) +
\gamma_\theta(\|\theta-\theta^\ast\|) + \gamma_\eta(\eta)$.
\textbf{Jump:} at each fault $t_k^\star$ the post-redistribution
geometry term satisfies
$V(t_k^{\star,+}) \le V(t_k^{\star,-}) + \chi(\Delta_f)$
(\cref{lem:jump-continuity}). \textbf{Contraction:} composing
flow with jumps over one inter-fault dwell yields the
per-cycle bound
$V(t_{k+1}^-) \le \rho\,V(t_k^-) + c$ with
$\rho = \exp(-\alpha_{\min}\tau_{\mathrm{pend}}) \approx 0.005$
(\cref{lem:dwell-decay}). The H2 hypothesis
$\tau_{d,k} \ge \tau_{\mathrm{pend}}$ is structural: it is the
minimum window over which the feed-forward identity
$T_i^{\mathrm{ff}} = T_i$ contracts the post-jump overshoot
back below $V(t_k^-)$.}
\label{fig:lyapunov-sawtooth}
\end{figure}
